\chapter{Objetivo}

Aunque este trabajo no se desarrolla directamente dentro del marco del proyecto Retinar, mi investigación se alinea con su propósito general, buscando aprovechar tecnologías basadas en IA para mejorar el tamizaje y diagnóstico de la RD en contextos con recursos limitados.

Muchas de las estrategias para ampliar el acceso al diagnóstico dependen de redes de captura de imágenes operadas por técnicos no especialistas, ya sea en entornos de atención primaria o mediante equipos móviles. En este contexto, garantizar que las imágenes capturadas sean de calidad suficiente para permitir un diagnóstico adecuado resulta esencial. Imágenes de mala calidad no solo dificultan la evaluación clínica, sino que obligan al paciente a realizar nuevas capturas, lo que conlleva múltiples inconvenientes: molestias, costos adicionales, pérdida de tiempo en el que la enfermedad puede progresar sin intervención, y un aumento del riesgo de subdiagnóstico en poblaciones vulnerables.

El objetivo de esta tesis es brindar una solución computacional que clasifique de forma automática las imágenes de fondo de ojo en las categorías de buena y mala calidad. Esto permitirá contar con una herramienta que pueda informar al técnico en el momento propio de la captura si la foto cumple con los criterios de calidad sufientes para continuar el proceso, o si debe realizar una nueva captura, evitando errores de diagnóstico posteriores y permitiendo repetir el estudio antes de que el paciente abandone la consulta.
